% ID: SEL-D0014-BB-E03
% ESTADO: cotejada. Segunda pasada visual completa y borrador compilado revisados.
% Fuente: Selectividad/matematicas_aplicadas/matematicas_aplicadas_P_B.pdf
% SHA-256: 51601c4f5a6c1dfcd860f621a659f8b9c09f53ff4cd31a1fa1e827dbd8f06159
% Página: 1; bloque B; ejercicio original 3.

\begin{enumerate}[label=\alph*)]
  \item \textbf{(1.2 puntos)} Se considera la función $f(x)=ax^2+bx+3$.
  Calcule los valores $a$ y $b$, sabiendo que la gráfica de $f$ pasa por el
  punto $(2,3)$ y que la pendiente de la recta tangente a la gráfica de $f$
  en dicho punto es $m=-2$.
  \item \textbf{(1.3 puntos)} Represente gráficamente la función
  $g(x)=-x^2+6x-5$ y calcule el área comprendida entre la gráfica de la
  función $g$, el eje de abscisas y las rectas $x=2$ y $x=4$.
\end{enumerate}
